\section{Capítulo 13 --- Agentes analíticos, controle e observabilidade}

Os itens distinguem LLM, ferramenta, skill, memória e coordenação. Um agente recebe apenas as permissões necessárias, registra decisões e recusa ou solicita intervenção quando não dispõe de evidência autorizada. Dados, custos e situações foram elaborados para fins didáticos.

\subsection{Questões objetivas}

\ItemHeader{1}{Objetiva}{Distinguir responsabilidades agentic}{Relatório mensal recorrente}
\TextoBase
Um fluxo recebe a solicitação, consulta métricas, compara o mês anterior, produz narrativa e pede aprovação humana. A equipe chama de ``agente'' tanto o modelo quanto uma função SQL e um manual de procedimento.
O relatório alimenta decisões de estoque, e cada etapa deve possuir entrada, saída, permissão e condição de parada verificáveis.
\FonteDidatica
\Comando A descrição tecnicamente adequada é
\begin{enumerate}[label=\Alph*)]
\item o LLM é a ferramenta, o SQL é o agente e o manual é a memória episódica.
\item o agente coordena estado e decisões; SQL é ferramenta; o manual pode constituir skill ou playbook; o LLM apoia interpretação.
\item todo componente é um agente porque participa do mesmo fluxo.
\item a memória decide o próximo passo e o agente apenas armazena texto.
\item a skill executa diretamente qualquer ação e torna permissões desnecessárias.
\end{enumerate}

\ItemHeader{2}{Objetiva}{Calcular confiabilidade composta}{Pipeline de três etapas}
\TextoBase
Sob independência aproximada no conjunto de teste, o planejador escolhe a ferramenta correta em 96\% dos casos, a ferramenta retorna dados corretos em 98\% e o verificador aprova corretamente em 95\%. O sucesso exige as três etapas.
O fluxo só entrega resposta quando planejamento, consulta e validação têm sucesso; falhar em qualquer etapa invalida a execução completa.
\FonteDidatica
\Comando A taxa ponta a ponta estimada é
\begin{enumerate}[label=\Alph*)]
\item 89,4\%, pelo produto $0{,}96\times0{,}98\times0{,}95$.
\item 96,3\%, pela média aritmética das três taxas.
\item 93,0\%, subtraindo do melhor resultado o pior erro.
\item 99,99\%, porque verificadores compensam qualquer falha anterior.
\item 78,0\%, somando as três porcentagens de erro.
\end{enumerate}

\ItemHeader{3}{Objetiva}{Aplicar menor privilégio por ferramenta}{Agente financeiro}
\TextoBase
O agente precisa ler uma SPEC de receita e abrir um ticket quando detectar divergência. A credencial atual também permite apagar tabelas e aprovar pagamentos.
A ferramenta de pagamento produz efeito externo irreversível, enquanto leitura de métricas e simulação podem operar com permissões distintas.
\FonteDidatica
\Comando A configuração coerente é
\begin{enumerate}[label=\Alph*)]
\item manter a credencial ampla e instruir o modelo a usar somente duas ações.
\item separar ferramentas de leitura e criação de ticket, com escopos mínimos, validação de parâmetros e aprovação para efeitos relevantes.
\item ocultar os nomes das ações perigosas no prompt, mantendo-as acessíveis.
\item usar memória para registrar que apagar tabela é proibido, dispensando controle técnico.
\item permitir pagamento quando a confiança textual superar 90\%.
\end{enumerate}

\ItemHeader{4}{Objetiva}{Calcular orçamento de tentativas}{Retry de consulta}
\TextoBase
Uma tentativa custa R\$0,12 e leva 1,4 s. Há orçamento de R\$0,40 e 5 s por solicitação, incluindo 0,6 s fixos fora das tentativas. Não existe paralelismo.
Cada tentativa custa R\$0,12 e 1,4 segundo; há ainda 0,6 segundo fixo, orçamento de R\$0,40 e prazo máximo de cinco segundos.
\FonteDidatica
\Comando O número máximo de tentativas completas permitido pelos dois orçamentos é
\begin{enumerate}[label=\Alph*)]
\item 2, porque a terceira ultrapassa o orçamento de tempo.
\item 3, pois custa R\$0,36 e usa $0{,}6+3(1{,}4)=4{,}8$ s.
\item 4, porque o custo de R\$0,48 pode ser compensado pela latência.
\item 1, porque qualquer retry produz efeito duplicado.
\item 5, dividindo o limite de tempo pela parte fixa.
\end{enumerate}

\ItemHeader{5}{Objetiva}{Escolher memória apropriada}{Correção de preferência e regra}
\TextoBase
Um usuário prefere gráficos em barras; a política corporativa determina que valores inferiores a cinco sejam suprimidos. O protótipo grava ambas as informações como lembranças livres produzidas pelo modelo.
A preferência de formato pertence ao usuário, mas uma política revogada não pode persistir como memória nem ser reintroduzida pelo histórico.
\FonteDidatica
\Comando O redesenho adequado é
\begin{enumerate}[label=\Alph*)]
\item armazenar preferência e política juntas em memória conversacional sem versão.
\item tratar preferência como dado de perfil revogável e política como regra governada, versionada e aplicada fora da memória livre.
\item converter a política em mensagem do usuário para que possa ser substituída em cada conversa.
\item descartar a preferência e manter somente o histórico integral indefinidamente.
\item permitir que o agente escolha qual regra lembrar conforme a tarefa.
\end{enumerate}

\ItemHeader{6}{Objetiva}{Dimensionar revisão humana}{Fila de HITL}
\TextoBase
O agente processa 1.500 casos/dia. Vinte por cento vão para revisão humana; cada revisão leva 6 minutos. Cada analista dispõe de 360 minutos úteis por dia.
Cada revisão leva seis minutos, existem 1.500 casos diários e 20\% são encaminhados; cada analista dispõe de 360 minutos produtivos.
\FonteDidatica
\Comando O número mínimo de analistas para absorver a fila diária é
\begin{enumerate}[label=\Alph*)]
\item 3, porque há 300 revisões.
\item 4, pois cada analista revisa 60 casos e cinco seriam necessários apenas com folga.
\item 5, porque são 1.800 minutos de revisão e cada analista oferece 360.
\item 6, dividindo casos totais por minutos de uma revisão.
\item 25, considerando todos os casos como revisão obrigatória.
\end{enumerate}

\ItemHeader{7}{Objetiva}{Decidir entre agente único e multiagente}{Análise de vendas}
\TextoBase
Um agente único chama duas ferramentas determinísticas e produz resposta em quatro passos estáveis. A proposta multiagente cria planejador, pesquisador, crítico e redator, elevando custo em 140\% sem ganho medido de correção.
A equipe deve justificar coordenação adicional por decomposição real da tarefa, e não pelo pressuposto de que mais agentes sempre melhoram.
\FonteDidatica
\Comando A decisão baseada em evidência é
\begin{enumerate}[label=\Alph*)]
\item adotar multiagente, porque mais papéis implicam maior inteligência.
\item manter o fluxo simples como baseline e exigir ganho mensurável ou necessidade real de especialização antes de ampliar coordenação.
\item remover todas as ferramentas e deixar agentes discutirem os números.
\item escolher aleatoriamente para evitar viés arquitetural.
\item duplicar cada agente para reduzir o custo total.
\end{enumerate}

\ItemHeader{8}{Objetiva}{Interpretar traces ponta a ponta}{Falhas em cem execuções}
\TextoBase
Em 100 execuções, 92 criaram plano válido, 88 chamaram ferramenta autorizada, 82 produziram resposta sustentada e 78 passaram no verificador final.
Os contadores representam etapas encadeadas; mudar denominadores entre elas esconderia onde o agente perdeu consultas ou falhou na validação.
\FonteDidatica
\Comando As taxas condicionais após o plano, após a ferramenta e após a sustentação, e o sucesso total são aproximadamente
\begin{enumerate}[label=\Alph*)]
\item 95,7\%; 93,2\%; 95,1\%; 78\%.
\item 88\%; 82\%; 78\%; 95,1\%.
\item 92\%; 88\%; 82\%; 78\%.
\item 95,7\%; 82\%; 78\%; 93,2\%.
\item 4\%; 6\%; 4\%; 22\%.
\end{enumerate}

\ItemHeader{9}{Objetiva}{Aplicar recusa segura}{Pergunta sem evidência autorizada}
\TextoBase
O usuário pede previsão de demissões por pessoa. O agente tem apenas dados agregados autorizados, e a política proíbe inferência individual. Uma ferramenta de busca não encontra fonte adicional.
A resposta será usada em decisão de crédito; ausência de fonte autorizada exige abstenção útil, sem fabricar aproximação nem executar ferramenta privilegiada.
\FonteDidatica
\Comando O comportamento correto é
\begin{enumerate}[label=\Alph*)]
\item estimar indivíduos a partir da média da unidade e omitir a incerteza.
\item recusar a inferência individual, explicar a restrição e oferecer análise agregada autorizada ou encaminhamento humano.
\item pedir ao modelo que invente dados sintéticos com nomes reais.
\item executar várias buscas até alguma resposta parecer confiante.
\item guardar a pergunta na memória e responder depois sem nova autorização.
\end{enumerate}

\ItemHeader{10}{Objetiva}{Aplicar gate de implantação agentic}{Comparação de versões}
\TextoBase
O gate exige sucesso útil $\geq80\%$, ações indevidas=0, p95 $\leq8$ s e custo $\leq$R\$0,50.
\begin{center}\small
\begin{tabularx}{\textwidth}{@{}Yrrrr@{}}\toprule
\textbf{Versão} & \textbf{Sucesso} & \textbf{Ações indevidas} & \textbf{p95} & \textbf{Custo}\\\midrule
Único & 81\% & 0 & 6,8 s & R\$0,31\\
Multi A & 87\% & 1 & 7,2 s & R\$0,48\\
Multi B & 84\% & 0 & 9,1 s & R\$0,46\\
Multi C & 82\% & 0 & 7,7 s & R\$0,56\\\bottomrule
\end{tabularx}\end{center}
\FonteDidatica
Os gates mínimos são zero ação indevida, correção útil de 80\%, p95 de sete segundos e custo máximo de R\$0,35.
\Comando A versão aprovada é
\begin{enumerate}[label=\Alph*)]
\item Multi A, por maximizar sucesso útil.
\item Multi B, por não realizar ações indevidas.
\item Multi C, por equilibrar sucesso e latência.
\item Único, por atender simultaneamente a todos os gates.
\item a média das versões multiagente, porque reduz violações individuais.
\end{enumerate}

\subsection{Questões discursivas}

\ItemHeader{11}{Discursiva}{Desenhar agente analítico governado}{Investigação de queda de receita}
\TextoBase
O agente deve consultar métricas certificadas, decompor receita, buscar incidentes e redigir hipótese. Não pode alterar dados nem atribuir causa sem evidência; conclusões de alto impacto precisam de revisão.
A investigação envolve métricas financeiras e pode acionar consulta, cálculo e visualização, mas nenhuma ferramenta deve publicar decisão sem validação.
\FonteDidatica
\Comando Desenhe estado, nós, ferramentas, skill, memória, gates, recusa, HITL, trace e critérios de parada. Explique como distinguir hipótese, evidência e decisão.

\ItemHeader{12}{Discursiva}{Calcular custo e benefício do multiagente}{Teste A/B arquitetural}
\TextoBase
Em 2.000 casos, o agente único obteve 1.600 respostas úteis a R\$0,28 por caso. O multiagente obteve 1.720 respostas úteis a R\$0,46 por caso. Ambos tiveram zero ação indevida. O orçamento incremental máximo é R\$4 por resposta útil adicional.
O teste contém duas mil solicitações: agente único custa R\$560 e resolve 1.600; multiagente custa R\$920 e resolve 1.720.
\FonteDidatica
\Comando Calcule custo total, custo por resposta útil, ganho de respostas e custo incremental por resposta adicional. Decida se o multiagente atende ao limite e indique outra evidência necessária.

\ItemHeader{13}{Discursiva}{Projetar política de memória}{Preferências, fatos e políticas}
\TextoBase
O agente mistura histórico, preferências, métricas calculadas e políticas em um vetor sem origem, validade ou opção de exclusão. Uma regra revogada continua influenciando respostas.
O projeto deve distinguir preferência revogável, fato com proveniência e política autoritativa, definindo retenção e precedência para cada classe.
\FonteDidatica
\Comando Proponha classes de memória, origem, versão, validade, autorização, retenção, correção/exclusão, recuperação e testes que impeçam política revogada e dado de outro usuário no contexto.

\ItemHeader{14}{Discursiva}{Calcular indicadores de HITL}{Piloto de decisões}
\TextoBase
Em 800 casos, 240 foram encaminhados a humanos. Revisores alteraram 72 decisões; 18 alterações ocorreram após o SLA. Entre os 560 não revisados, auditoria amostral de 140 encontrou 7 erros relevantes.
Uma amostra aleatória de 140 casos não encaminhados revelou sete erros, informação necessária para avaliar segurança fora da fila humana.
\FonteDidatica
\Comando Calcule taxa de encaminhamento, taxa de alteração entre revisados, atraso entre alterações e erro estimado nos não revisados. Interprete se encaminhar menos casos é evidência suficiente de melhora e proponha ajuste verificável.

\ItemHeader{15}{Discursiva}{Responder a incidente agentic}{Ferramenta realizou efeito indevido}
\TextoBase
Um agente abriu 83 tickets duplicados após timeout porque repetiu a ferramenta sem chave idempotente. Traces guardam argumentos e IDs, mas também contêm e-mails em claro.
A ferramenta enviou mensagem externa sem confirmação; a análise deve identificar escopo, credencial, trace e usuários potencialmente afetados.
\FonteDidatica
\Comando Proponha contenção, correção dos tickets, preservação minimizada de evidência, comunicação, idempotência, orçamento de retry, confirmação de efeito, testes e critérios de reativação.

\newpage
\subsection{Respostas explicadas das questões pares}

\paragraph{Questão 2 — resposta A.} O sucesso composto é $0{,}96\times0{,}98\times0{,}95=0{,}89376$, ou 89,4\%. A média das etapas superestima o fluxo que exige sucesso simultâneo.
\[P(\text{sucesso})=0{,}96(0{,}98)(0{,}95)=0{,}89376\]

\paragraph{Questão 4 — resposta B.} Três tentativas custam R\$0,36 e usam 4,8 s com a parte fixa. Quatro custariam R\$0,48 e usariam 6,2 s, violando ambos os orçamentos.

\paragraph{Questão 6 — resposta C.} São $1.500(0{,}20)=300$ revisões, ou 1.800 minutos. $1.800/360=5$ analistas, sem margem adicional.

\paragraph{Questão 8 — resposta A.} $88/92\approx95{,}7\%$; $82/88\approx93{,}2\%$; $78/82\approx95{,}1\%$; sucesso total $78/100=78\%$. O funil localiza perdas sem mudar denominadores silenciosamente.

\paragraph{Questão 10 — resposta D.} O agente único alcança 81\%, zero ação indevida, 6,8 s e R\$0,31. Cada versão multiagente viola segurança, latência ou custo; maior correção não compensa um gate obrigatório.

\paragraph{Questão 12 — critérios.} Único: R\$560 total e R\$0,35 por útil. Multiagente: R\$920 total e aproximadamente R\$0,535 por útil. O ganho é 120 respostas por R\$360 incrementais, ou R\$3 por resposta adicional; passa o limite de R\$4. A decisão ainda requer latência, distribuição de erros, carga operacional e avaliação por subgrupo.

\paragraph{Questão 14 — critérios.} Encaminhamento $=240/800=30\%$; alteração $=72/240=30\%$; atraso $=18/72=25\%$; erro amostral nos não revisados $=7/140=5\%$. Menor encaminhamento não prova melhora: pode aumentar erro não detectado. Ajustes válidos usam amostra auditada, risco calibrado, capacidade e SLA.
